% Self-contained contribution; use the master theorem counters when included.
% All statements below are unconditional and proved.  No conclusion about
% existence or nonexistence of an ordinary all-depth ghost is asserted.
\subsection{Integral lifts of one fixed pair form a compact tower}
\label{sec:ghost-compact-tower}

Let \(G=G_{\mathbb Q}\), let \(F_\ell\) be the free pro-\(\ell\)
group on the two geometric meridians, and set
\[
 N_{r,\ell}=F_\ell/\Gamma_{r+1}F_\ell,\qquad
 L_{r,\ell}=\Gamma_rF_\ell/\Gamma_{r+1}F_\ell.
\]
All commutator subgroups here are closed.  The chosen rational
tangential section gives the actual continuous action of \(G\).
Its action on the abelianization is scalar multiplication by the
\(\ell\)-adic cyclotomic character \(\chi_\ell\).  Consequently
\[
       L_{r,\ell}\simeq\mathbb Z_\ell(r)^{b_r},
       \tag{*}\label{eq:ghost-central-layer}
\]
where \(b_r\) is the number of basic commutators of weight \(r\).
Indeed commutators induce multilinear maps on the associated graded
group, its Hall basis is generated in weight one, and applying the
scalar \(\chi_\ell(g)\) to each of the \(r\) entries multiplies
every weight-\(r\) commutator by \(\chi_\ell(g)^r\).

Fix one \emph{literal} continuous cocycle
\(a:G\to N_{1,\ell}=\mathbb Z_\ell(1)^2\).  Write
\[
 \mathscr S_{r,\ell}(a)=
 \{c\in Z^1_{\mathrm{cts}}(G,N_{r,\ell}):c\bmod\Gamma_2=a\}.
\]
We put the compact-open topology on these spaces of continuous maps.
They consist of actual cocycles, not conjugacy classes.

\begin{lemma}[Compact higher-weight cocycles]
\label{lem:ghost-higher-z1}
For every \(r\geq2\), the module
\(Z^1_{\mathrm{cts}}(G,\mathbb Z_\ell(r))\) is a finite free
\(\mathbb Z_\ell\)-module, and its module topology is its
compact-open topology.  Every such cocycle factors through
\(G_{\mathbb Q,\{\ell,\infty\}}\).
\end{lemma}
\begin{proof}
Put \(T=\mathbb Z_\ell(r)\). At a prime \(q\ne\ell\), inertia
acts trivially on \(T\), so the restriction of a cocycle to inertia
is a continuous homomorphism. Wild inertia and the prime-to-\(\ell\)
part of tame inertia have zero image in the additive pro-\(\ell\)
group \(T\). For a tame \(\ell\)-generator \(\tau\) and an
arithmetic Frobenius \(\sigma\), the relation
\(\sigma\tau\sigma^{-1}=\tau^q\), after wild inertia has been
removed, and the cocycle identity give
\[
        q\,c(\tau)=c(\sigma\tau\sigma^{-1})
                    =q^r c(\tau).
\]
The equality \(q^r-q\ne0\) and torsion-freeness of \(T\) give
\(c(\tau)=0\). The action on \(T\) is unramified outside
\(\ell\). The cocycle therefore vanishes on the closed normal
subgroup generated by these inertia groups and descends to
\(G_S=G_{\mathbb Q,\{\ell,\infty\}}\).

For \(n\geq1\), the set \(Z^1(G_S,T/\ell^nT)\) is finite.
To see this without a finite-generation assertion about \(G\),
send a cocycle \(c\) to the homomorphism
\[
 g\longmapsto(c(g),\chi_\ell(g)^r\bmod\ell^n)
 \quad\hbox{into}\quad
 (T/\ell^nT)\rtimes\operatorname{im}(\chi_\ell^r\bmod\ell^n).
\]
Its kernel cuts out an extension of bounded degree unramified
outside the fixed finite set \(S\). There are only finitely many
such extensions by Hermite--Minkowski with fixed ramification;
each finite quotient admits only finitely many maps to this target.
This proves finiteness. The arithmetic finiteness assertion used
here is exactly \cite[Theorem~8.42]{GhostMilneANT}.

Compatible reductions give the topological isomorphism
\[
 Z:=Z^1(G_S,T)=\varprojlim_n Z^1(G_S,T/\ell^nT).
\]
In fact a compatible sequence defines a continuous map to \(T\)
and the cocycle identities hold in every reduction. Conversely,
each continuous map has these reductions. Uniform congruence
modulo \(\ell^n\) gives the compact-open topology because the
source is compact. Thus \(Z\) is compact Hausdorff.

Reduction modulo \(\ell\) has kernel precisely \(\ell Z\): a
cocycle with values in \(\ell T\) may be divided continuously
by \(\ell\), and its cocycle identity divides as well. Hence
\(Z/\ell Z\) is finite. Choose lifts \(z_1,\ldots,z_d\) of a
basis of this finite vector space. Repeatedly reducing a residual
element modulo \(\ell Z\) expresses any \(z\in Z\), modulo
\(\ell^nZ\), as \(\sum_i a_{i,n}z_i\), with the coefficients
compatible modulo \(\ell^n\). Their limits \(a_i\in\mathbb Z_\ell\)
exist. The remaining cocycles have values in \(\ell^nT\), so
converge uniformly to zero. It follows that
\(z=\sum_i a_i z_i\). This proves finite generation algebraically,
not merely the finiteness of a collection of finite quotients.
The module is torsion-free as a submodule of the continuous
\(T\)-valued functions, and a finitely generated torsion-free
module over \(\mathbb Z_\ell\) is free. Its module topology
agrees with the topology above because its basis gives a continuous
bijection from a compact power of \(\mathbb Z_\ell\) to a
Hausdorff space.
\end{proof}

\begin{lemma}[The integral Hall coordinate model]
\label{lem:ghost-hall-model}
For fixed \(r\), ordered basic commutators of weights at most
\(r\) identify \(N_{r,\ell}\), as a topological space, with
\(\mathbb Z_\ell^{h_r}\). Multiplication, inversion, and taking
\(\mathbb Z_\ell\)-powers have coordinates given by rational
polynomials. Their denominators have a fixed bound depending on
\(r\). The quotient to \(N_{r-1,\ell}\) deletes the weight-\(r\)
coordinates, and setting those coordinates equal to zero gives a
continuous section of this quotient as a map of spaces.

The coordinate functions of \(g(u)\), for the actual action of
\(G\), are rational polynomials in the coordinates of \(u\),
with coefficients that are continuous functions of \(g\).
They too have bounded denominators at each fixed \(r\).
\end{lemma}
\begin{proof}
Start with the free discrete nilpotent group of class \(r\),
using a Lyndon Hall ordering of the basic commutators.
The collecting process orders a word by repeatedly using
\(vu=uv[v,u]\); the new commutator has larger weight, so the
process terminates after weight \(r\). At each weight the basic
commutators form the integral basis obtained by imposing
antisymmetry and the Jacobi relation on brackets in the two
generators. Their independence follows by expanding into the
free associative algebra: the successive leading Hall words
have coefficient one and distinct leading words. Removing
these successively proves that every element has a unique
expression \(\prod_i b_i^{a_i}\), with \(a_i\in\mathbb Z\).
The same collection, by induction on weight and first for
nonnegative exponents, gives rational polynomial formulas for
products and powers. At the next weight the only new terms
are finite iterated sums of the already obtained polynomial
terms; discrete summation of a polynomial introduces only
binomial polynomials. The number of weights and commutator
types is finite at fixed class. The resulting polynomial
identities extend to negative integer exponents by the same
collection identities. Inversion follows by collecting the reversed
word with negated exponents. Thus there are only finitely many
denominators at fixed class; no claim that these denominators are
units in \(\mathbb Z_\ell\) is required. The explicit
polynomial-recursion implementation of this integral collecting
argument is given in \cite[\S\S3--4]{GhostCantEick2019};
the completion argument below supplies the required comparison
with the free pro-\(\ell\) quotient.

Extend these rational polynomials from \(\mathbb Z^{h_r}\) to
\(\mathbb Z_\ell^{h_r}\). Their values are integral by density,
since the original values are integral and \(\mathbb Z_\ell\)
is closed in \(\mathbb Q_\ell\). All group identities extend
by continuity. We obtain a compact totally disconnected group
\(H_r\), with the discrete free nilpotent group dense in it.
It is a pro-\(\ell\) group. Indeed each one-parameter subgroup
\(a\mapsto b_i^a\) is a continuous image of \(\mathbb Z_\ell\).
Any finite quotient of \(H_r\) is nilpotent and is generated
by the images of the \(b_i\), each of which has \(\ell\)-power
order. Its Sylow decomposition therefore has only the
\(\ell\)-factor.

We also check the exact universal completion, including small
\(\ell\). A homomorphism from the discrete group to a finite
\(\ell\)-group of exponent \(\ell^e\) kills every collected
word all of whose coordinates belong to \(\ell^e\mathbb Z\).
This coordinate box is a neighbourhood of the identity in the
induced topology from \(H_r\), whether or not the box is itself
a subgroup. Thus the homomorphism is uniformly continuous and
extends uniquely to \(H_r\). Every finite \(\ell\)-group of
class at most \(r\) is included in this argument. It proves that
\(H_r\) is exactly the pro-\(\ell\) completion of the discrete
free nilpotent group, and the universal property of the free
pro-\(\ell\) group identifies it with \(N_{r,\ell}\).

Finally an endomorphism of this group is determined by the
images of its two free generators. Substitute these images into
the finitely many basic commutator words, then use the product
and power polynomials to evaluate a collected word. This gives
rational polynomials jointly in the input coordinates and the
coordinates of the two images, with fixed denominators. The
images of the two generators under \(g\) are continuous functions
of \(g\), taking values in the compact integral coordinate space.
Substitution gives precisely the assertion about the actual
arithmetic action, including its bounded denominators.
\end{proof}

\begin{lemma}[Finite polynomial families of integral primitives]
\label{lem:ghost-primitive-family}
Let \(T\) be a finite free \(\mathbb Z_\ell\)-module with a
continuous \(G\)-action. Fix continuous rational two-cochains
\(c_1,\ldots,c_m\in C^2(G,T\otimes\mathbb Q_\ell)\). Then
\[
 R=\left\{u\in\mathbb Z_\ell^m:
       \sum_i u_i c_i\in\delta C^1(G,T)\right\}
\]
is a closed finite free submodule of \(\mathbb Z_\ell^m\).
There is a continuous \(\mathbb Z_\ell\)-linear map
\(h:R\to C^1(G,T)\) satisfying
\(\delta h(u)=\sum_i u_i c_i\). Each coordinate of \(h\),
relative to finitely many fixed continuous cochains, is the
restriction to \(R\) of a \(\mathbb Q_\ell\)-linear form on
\(\mathbb Q_\ell^m\).
\end{lemma}
\begin{proof}
The displayed set is a submodule of the finite free module
\(\mathbb Z_\ell^m\). The ring \(\mathbb Z_\ell\) is a
Noetherian principal ideal domain, so \(R\) has a finite basis
\(r_1,\ldots,r_k\). The matrix with these columns defines a
continuous injection \(\mathbb Z_\ell^k\to\mathbb Z_\ell^m\).
Its image is compact and therefore closed; this image is exactly
\(R\). For each basis vector, choose an integral continuous
cochain \(h_j\) with
\(\delta h_j=\sum_i(r_j)_i c_i\), as the definition of \(R\)
permits. The matrix has a \(\mathbb Q_\ell\)-linear left inverse
on its column span; extend it linearly to \(\mathbb Q_\ell^m\).
Its component forms \(\lambda_j\) are integral on \(R\), and
\[
                    h(u)=\sum_j\lambda_j(u)h_j
\]
has all asserted properties. Uniform continuity on the compact
source of the cochains gives continuity into the compact-open
function space. This proof does not assert that the full coboundary
module \(\delta C^1(G,T)\) is closed in the full cochain space.
\end{proof}

\begin{proposition}[Compactness for a fixed abelian cocycle]
\label{prop:ghost-compact-fibres}
Every \(\mathscr S_{r,\ell}(a)\) is compact Hausdorff. More
precisely, it is the continuous image of a closed subset
\(X_r\subseteq\mathbb Z_\ell^{d_r}\), and its Hall coordinate
functions are rational polynomials in these parameters, with
continuous cochain coefficients.
\end{proposition}
\begin{proof}
We prove the stronger parametrization assertion by induction.
At \(r=1\), take a one-point parameter space and the fixed
cocycle \(a\). Suppose \(c_t\), \(t\in X_r\), parametrizes
all cocycles at level \(r\). Lift each using the Hall section
that sets its new central coordinates to zero. The defect
\[
 C_t(g,h)=\widetilde c_t(g)\,
          g(\widetilde c_t(h))\,
          \widetilde c_t(gh)^{-1}
\]
lies in \(L_{r+1,\ell}\) for \(t\in X_r\). By
Lemma~\ref{lem:ghost-hall-model} its central coordinates are a
finite rational polynomial in \(t\), so they have an expression
\[
                  C_t=\sum_{i=1}^m v_i(t)c_i,
\]
where the \(v_i\) are integral monomials in the parameters and
the \(c_i\) are fixed continuous rational two-cochains.
The individual \(c_i\) need not be cocycles. On \(X_r\), the
sum is the integral central defect, which is a two-cocycle.
Apply Lemma~\ref{lem:ghost-primitive-family} to these cochains.
The exact locus on which the lifted cocycle can be corrected is
\[
          X'_r=\{t\in X_r:(v_1(t),\ldots,v_m(t))\in R\}.
\]
This is closed and compact. Multiplying the Hall lift by the
central cochain \(-h(v(t))\) kills its defect. All further
corrections are exactly the literal cocycles
\(Z^1(G,L_{r+1,\ell})\), which form a finite free compact module
by \((\ast)\) and Lemma~\ref{lem:ghost-higher-z1}. Choosing a
basis gives the parameter space
\(X_{r+1}=X'_r\times\mathbb Z_\ell^{e_{r+1}}\).
Every lift is represented because the difference of two central
corrections with the same defect is such a literal one-cocycle.

The linear forms in Lemma~\ref{lem:ghost-primitive-family},
composed with the monomials \(v_i(t)\), are rational polynomial
functions. Their values are integral on \(X'_r\). Adding the
new cocycle parameters and applying the Hall product formulas
preserves the claimed polynomial parametrization. This completes
the induction. The joint evaluation map
\(X_r\times G\to N_{r,\ell}\) is continuous by these explicit
formulas, hence its adjoint to the compact-open function space
is continuous. Its image is exactly
\(\mathscr S_{r,\ell}(a)\), proving compactness. The function
space is Hausdorff because its target is Hausdorff.
\end{proof}

\begin{theorem}[Exact compatibility equivalence for one fixed pair]
\label{thm:ghost-fixed-pair-compatible}
For any fixed power-completed Kummer pair \((A,B)\), the following
two statements are equivalent:
\begin{enumerate}
 \item For every integer \(r\geq1\), there exists a continuous
 section of \(E/\Gamma_{r+1}\Pi\to G_{\mathbb Q}\) whose
 abelian Kummer class is \((A,B)\).
 \item There exists a continuous section of
 \(E/\bigcap_r\Gamma_{r+1}\Pi\to G_{\mathbb Q}\) whose
 abelian Kummer class is \((A,B)\).
\end{enumerate}
The corresponding equivalence also holds with the requirement
that all restrictions be geometrically conjugate to one fixed
family of local point sections, or of local cuspidal sections,
at every level.
\end{theorem}
\begin{proof}
Fix \(\ell\) and choose a literal abelian cocycle \(a\)
representing its component of \((A,B)\). Any lift with this
abelian cohomology class can be geometrically conjugated to
have literal abelian cocycle \(a\): their difference is the
coboundary of an element of \(N_{1,\ell}\), and that element
has a preimage in \(N_{r,\ell}\). Thus existence at every
level means every \(\mathscr S_{r,\ell}(a)\) is nonempty.
The transition maps are continuous. In the compact product
\(\prod_r\mathscr S_{r,\ell}(a)\), the equations saying that
the adjacent coordinates project to each other are closed.
Every finite collection of these equations has a solution:
choose a section at the largest occurring level and project it
to the lower levels involved, and choose arbitrary elements at
the other coordinates. Compactness gives a simultaneous
solution to all equations. Its coordinate maps define a
continuous section of \(F_\ell\rtimes G\to G\), since
\(F_\ell=\varprojlim_rN_{r,\ell}\). This last equality follows
also from the fact that every finite \(\ell\)-group is nilpotent.

The maximal pronilpotent geometric quotient is exactly
\(\prod_\ell F_\ell\): its finite quotients are precisely the
finite nilpotent quotients, whose Sylow factors are finite
\(\ell\)-quotients. The actual arithmetic action preserves
each characteristic factor. Taking the constructed cocycles
over all \(\ell\) therefore gives a continuous cocycle into
this product and the required section. Conversely a section
into the inverse limit projects to every level.

For the final assertion, fix a decomposition embedding at each
place. At level \(r\), the conjugacy orbit of the chosen local
cocycle is the continuous image of the compact group
\(N_{r,\ell}\) in its local compact-open function space. It is
therefore compact and closed. Restriction of global cocycles
is continuous, so imposing any collection of these local
matching conditions cuts out a closed subset of the compact
global section space. Projection preserves the conditions by
projecting the actual conjugators. The same compact-product
argument gives compatible globally defined sections at every
level satisfying them. For each place, the sets of geometric
conjugators at the successive levels are nonempty compact
sets with continuous transition maps. Their inverse limit is
nonempty by the same finite-intersection argument, giving an
actual conjugator in \(F_\ell\). Finally the product of these
conjugators over \(\ell\) gives the conjugator in the full
pronilpotent geometric group.
\end{proof}

This theorem keeps the pair \((A,B)\) fixed throughout. The
existence of a different pair at each depth does not supply its
first statement. It also begins with integral continuous lifts
at each depth. Replacing these by unrelated lifts in every
finite coefficient quotient is not one of its premises or
conclusions: the exact spaces above use continuous
\(\mathbb Z_\ell\)-valued cochains at each fixed class.
For the ordinary pairs of Part C, finite-support reconstruction
is already proved; the remaining all-depth question is whether
the equivalent sets in this theorem are empty for every
ordinary infinite-support pair.
